

\textbf{Inferring stochastic dynamical systems from snapshots with structured domain knowledge using Schrödinger Bridges.} Systems biology forms one example where data is scarce and noisy, but substantial domain knowledge exists. For example, practitioners often aim to infer an unobserved population trajectory using sample snapshots at multiple time points. E.g. given single-cell sequencing data, scientists would like to learn how gene expression changes over a cell’s life cycle. But sequencing any cell destroys that cell. 
So we can access data for any particular cell only at a single time point, but we have data across many cells. The deep learning community has recently explored using Schrödinger bridges (SBs) and their extensions in similar settings. However, existing methods either (1) interpolate between just two time points or (2) require a single fixed reference dynamic (often set to Brownian motion within SBs). But learning piecewise from adjacent time points can fail to capture long-term dependencies. And practitioners are typically able to specify a model family for the reference dynamic but not the exact values of the parameters within it. So we propose a new method that (1) learns the unobserved trajectories from sample snapshots across multiple time points and (2) requires specification only of a family of reference dynamics, not a single fixed one. We demonstrate the advantages of our method on simulated and real data. This work, \cite{shen2025multi}, was published as a conference paper on AISTATS 2025 with me as a co-first author. 


\textbf{Forecasting stochastic dynamical systems from snapshots with unknown volatility and missing observations using least square method in distributional space.} %Scientists often need to learn the underlying stochastic dynamics of systems from population-level snapshot data, where individual trajectories are unavailable. For example, in timecourse single-cell mRNA profiling, cellular transcriptional states measurement is performed on different biological replicates, because the measurement process destroys the cells.  
In the work above, I extended  existing Schrödinger bridge techniques to allow a \textit{family} of references in order to incoperate structural knowledge rather than using a fixed reference. However, existing methods based on Schrödinger bridge techniques, my own work included, rely on Kullback–Leibler (KL) divergence and assume known, constant volatility, limiting their applicability in realistic settings where volatility may be unknown or varying. In this follow up work, I propose a new framework that directly matches the joint distribution of the state (e.g., mRNA expression levels) and the time of observation using maximum mean discrepancy. This approach reduces to a least-squares formulation in distributional space and motivates an $R^2$-type goodness-of-fit measure for model inspection and comparison. We show in our experiments that the proposed method outperforms existing Schrödinger bridge–based baselines in forecasting and is robust to unknown volatility and missing observations. The work, \cite{berlinghieri2025oh}, is currently under submission with me as a co-first author.


%In biology, researchers want to understand how gene expression changes over a cell’s life cycle, and how disruptions of these processes can lead to cancer. The challenge is that single-cell sequencing, the way to measure such gene expression at cell level destroys each cell in the process, revealing only static ``snapshots'' of cells, not continuous trajectories. This makes trajectory inference an ill-posed inverse problem: classic time-series methods fail when there are no time series.

%Recent deep learning work has turned to Schrödinger bridges, extensions of diffusion generative models. However, existing approaches impose strong assumptions: many interpolate only between two time points, or assume a fixed reference process such as Brownian motion, or rely on knowing the system’s volatility, or assumed absence of missing observations. These assumptions break down in biology, where long-term dependencies matter, domain knowledge provide only structural references, and noise levels are unknown.

%In \cite{shen2025multi}, we developed methods that relax these restrictions while retaining tractability. First, we introduced multi-marginal Schrödinger bridges that learn trajectories across many time points, requiring only a family of plausible reference dynamics instead of a fixed one. This allowed us to capture realistic dynamics in biological clocks, predator–prey systems, ocean currents, and single-cell sequencing where prior approaches failed. This work was selected as a selective oral presentation at AISTATS 2025.

%Second, we replaced the Kullback–Leibler divergence, which is fragile to mismatched volatility or time horizons, with a maximum mean discrepancy (MMD) loss. This extends least-squares estimation to distributions and enables trajectory inference even when noise structure is unknown, enabled forecasting and handling of missing observations. On both synthetic and real datasets — from biological clocks and immune cell activation to predator–prey systems — this approach consistently outperformed baselines, offering more reliable forecasting and interpolations.

%Together, these contributions show how embedding flexible domain knowledge into probabilistic machine learning can resolve ill-posedness without requiring unrealistic assumptions. Looking forward, this line of work provides the foundation for my future research on multimodal stochastic dynamics, where we aim to extend Schrödinger bridges to handle Perturb-seq data and other settings with multiple, interacting observational sources.


%Ill-posed inverse problems are common in systems biology, but substantial domain knowledge exists. For example, biologists seek to understand how gene expression changes during cells' life cycle, and how malfunctions in this process lead to cancer. Although single-cell sequencing reveals gene expression, it destroys the cell, leaving only snapshots instead of full trajectories. This lack of trajectories makes classic time series methods ineffective and renders inference of the underlying dynamics ill-posed. To address the ill-posedness, we developed methods that incorporate domain knowledge—constrained enough to resolve ill-posedness, yet flexible enough to allow imperfect knowledge.

%To solve the trajectory inference problem, recent deep learning work has turned to Schrödinger bridges, extensions of diffusion generative models. However, Existing methods either interpolate locally or assume fixed, known references and noise, which fail in realistic biological systems. In \cite{shen2025multi}, we developed a method that (1) learns unobserved trajectories from snapshots across multiple time points and (2) requires only a family of reference dynamics rather than a fixed one. In experiments on biological clocks, ocean currents, single-cell sequencing, and predator–prey systems, this approach successfully captured dynamics where existing methods failed. This work was accepted as a selective oral presentation at AISTATS 2025. In an early version \cite{shen2024learning} we also highlighted examples of nonidentifiability from snapshot data that are used in many follow up works on identifiability.% and plan to develop a general identifiability theory for trajectory inference.  

%Existing Schrödinger bridge methods also assume that the system to have known noise structure, fixed time horizons and no missing observations. While this assumption simplifies optimization via Kullback–Leibler divergence, these assumptions makes it difficult to learn systems with unknown noise structure,  undermines forecasting and cannot handle missing observations (e.g., missing protein measurements), a key requirement in biological applications. In \cite{berlinghieri2025oh}, we proposed replaced the Kullback–Leibler divergence with a maximum mean discrepancy (MMD) loss. This extends least-squares estimation to distributions and enables trajectory inference even when noise structure is unknown and potentially some dimensions are not observed. On both synthetic and real datasets — from biological clocks and immune cell activation to predator–prey systems — this approach consistently outperformed baselines, offering more reliable forecasting and can handle missing observations. This work is currently under submission. 